\subsection{Functional calculus in a non-unital algebra}

Let A be a stellar algebra and let \(\widetilde{\mathbf{A}}\) be the unital stellar algebra deduced from A by adjoining a unit element \(e\) . Denote by \(\pi\) the hermitian character \(x + \lambda e \mapsto \lambda\) of \(\widetilde{\mathbf{A}}\) into \(\mathbf{C}\) ; we have \(\operatorname{Ker}(\pi) = \mathbf{A}\) .

Let \(x \in \mathrm{A}\) be a normal element. It is normal in \(\widetilde{\mathrm{A}}\) and \(\mathrm{Sp}_{\widetilde{\mathrm{A}}}^{\prime}(x) = \mathrm{Sp}_{\mathrm{A}}^{\prime}(x)\) . Denote by \(\mathcal{C}'(\mathrm{Sp}_{\mathrm{A}}^{\prime}(x))\) the stellar algebra of continuous functions \(f\) on \(\mathrm{Sp}_{\mathrm{A}}^{\prime}(x)\) such that \(f(0) = 0\) .

Let \(f \in \mathcal{C}(\mathrm{Sp}_{\mathrm{A}}^{\prime}(x))\) . Since \(\pi(f(x)) = f(\pi(x))\) (prop. 8 of I, p.~112), we have \(f(x) \in \mathrm{A}\) if and only if \(f(0) = 0\) . The map \(f \mapsto f(x)\) defines a morphism of involutive algebras from the stellar algebra \(\mathcal{C}'(\mathrm{Sp}_{\mathrm{A}}^{\prime}(x))\) to A for which the image of the identity map \(z\) of \(\mathrm{Sp}_{\mathrm{A}}^{\prime}(x)\) is \(x\) . This morphism is isometric and its image is the stellar subalgebra of A generated by \(x\) .

PROPOSITION 12. --- The map \(f \mapsto f(x)\) is the unique morphism of involutive algebras from the stellar algebra \(\mathcal{C}'(\mathrm{Sp}_{\mathrm{A}}'(x))\) to A such that the identity map \(z\) of \(\mathrm{Sp}_{\mathrm{A}}'(x)\) has image \(x\) .

The elements \(z\) and \(\overline{z}\) of \(\mathcal{C}'(\mathrm{Sp}_{\mathrm{A}}'(x))\) generate a dense subalgebra of \(\mathcal{C}'(\mathrm{Sp}_{\mathrm{A}}'(x))\) (cf.~TG, X, p.~40, cor. 2). Since every morphism of involutive algebras from the stellar algebra \(\mathcal{C}'(\mathrm{Sp}_{\mathrm{A}}'(x))\) to the stellar algebra A is continuous (I, p.~104, prop. 2), the result follows.

The results of the preceding number concerning functional calculus extend to the general case. We shall simply state them and leave it to the readers to complete the proofs, mutatis mutandis.

PROPOSITION 13. --- One has the following properties:

\begin{enumerate}
\def\labelenumi{\alph{enumi})}
\item
  For every \(f \in \mathcal{C}'(\mathrm{Sp}_{\mathrm{A}}'(x))\) , one has \(\mathrm{Sp}_{\mathrm{A}}'(f(x)) = f(\mathrm{Sp}_{\mathrm{A}}'(x))\) ;
\item
  For every \(f \in \mathcal{C}'(\mathrm{Sp}_{\mathrm{A}}'(x))\) and for every \(g \in \mathcal{C}'(\mathrm{Sp}_{\mathrm{A}}'(f(x)))\) , one has \((g \circ f)(x) = g(f(x))\) ;
\item
  Let A' be a stellar algebra and let \(\pi\) be a morphism from A to A'; then \(\pi(x)\) is normal in A', one has \(\mathrm{Sp}_{\mathrm{A}'}^{\prime}(\pi(x)) \subset \mathrm{Sp}_{\mathrm{A}}^{\prime}(x)\) and \(\pi(f(x)) = f(\pi(x))\) for every \(f \in \mathcal{C}'(\mathrm{Sp}_{\mathrm{A}}^{\prime}(x))\);
\item
  If A is commutative, and if \(f \in \mathcal{C}'(\mathrm{Sp}_{\mathrm{A}}'(x))\) , then \(\mathcal{G}_{\mathrm{A}}'(f(x)) = f \circ \mathcal{G}_{\mathrm{A}}'(x)\) .
\end{enumerate}

Remark. --- Let A be a unital stellar algebra and let \(\widetilde{\mathrm{A}}\) be the unital stellar algebra deduced from A by adjoining a unit element \(e\) . For every \(x \in \mathrm{A}\) , one has \(\mathrm{Sp}_{\mathrm{A}}^{\prime}(x) = \mathrm{Sp}_{\mathrm{A}}(x) \cup \{0\}\) . Let \(x\) be a normal element of A. It is then a normal element of \(\widetilde{\mathrm{A}}\) and one therefore has two functional calculus maps in A, the first defined on \(\mathcal{C}(\mathrm{Sp}_{\mathrm{A}}(x))\) and the second on \(\mathcal{C}'(\mathrm{Sp}_{\mathrm{A}}^{\prime}(x))\) . Let \(f' \in \mathcal{C}'(\mathrm{Sp}_{\mathrm{A}}^{\prime}(x))\) ; if one denotes by \(f\) its restriction to \(\mathrm{Sp}_{\mathrm{A}}(x)\) , we then have \(f'(x) = f(x)\) .

